% Slide — PSO antes/depois
\begin{frame}{Particle Swarm Optimization -- \antes\ vs. \depois}
  \begin{columns}[T,onlytextwidth]
    \column{0.48\textwidth}
      \antes\\[0.3em]
      \small NumPy puro, laço Python por partícula
      \begin{itemize}\footnotesize
        \item 40 partículas $\times$ 80 iterações.
        \item $w=0{,}7$, $c_1=c_2=1{,}5$, \textbf{sem} fator de constrição.
        \item Velocidade sem limite explícito -- pode divergir.
        \item Avaliação: laço \texttt{for i in range(n\_particles)}, uma
          simulação por partícula.
      \end{itemize}
    \column{0.48\textwidth}
      \depois\\[0.3em]
      \small \texttt{TorchPSO}, avaliação em lote
      \begin{itemize}\footnotesize
        \item 40 $\times$ 80 (\textbf{inalterado}).
        \item Mesmos $w$, $c_1$, $c_2$ \highlight{+ fator de constrição de
          Clerc \& Kennedy (2002)}:
          $\chi = \dfrac{2}{|2-\phi-\sqrt{\phi^2-4\phi}|}$, $\phi=c_1{+}c_2$.
        \item $v \leftarrow \chi\,[w v + c_1 r_1(p_{\text{best}}{-}x) +
          c_2 r_2(g_{\text{best}}{-}x)]$ -- converge sem \textit{clamp}
          arbitrário de velocidade.
      \end{itemize}
  \end{columns}
  \vspace{0.8em}

  \begin{alertblock}{Sintoma observado que motivou a correção}
    \small Nos resultados anteriores, $Q$ frequentemente ``colava'' no
    limite superior do espaço de busca -- assinatura clássica de
    velocidade divergente em PSO sem constrição.
  \end{alertblock}

  \note{
    O PSO antigo usava a combinação w=0,7 e c1=c2=1,5 sem nenhum mecanismo
    de controle de velocidade -- combinação em que, pela teoria de
    convergência de PSO, a velocidade PODE divergir. O sintoma prático era
    Q colando repetidamente no teto do espaço de busca, o que é
    exatamente o que se espera quando a velocidade cresce sem controle: a
    partícula "voa" para fora da região de interesse e é sempre truncada
    no limite. O fator de constrição de Clerc e Kennedy resolve isso
    matematicamente, sem precisar de um clamp arbitrário de velocidade que
    seria mais um parâmetro para justificar.
  }
\end{frame}
